% =====================================================================================
%  Chapter 3 — System Modelling and Residual-Force Formulation
%  Thesis: Comparison of Control Strategies for Interaction-Force Aware Multirotor Teams
%
%  DRAFT, 23 August 2026 (writing-track item W.6).
%
%  SOURCES.
%  Claims about prior work: verify against the paper itself, never against a summary or ledger
%  in docs/ (see the ch2 header). This file is the source of truth for what it states.
%  Claims about OUR system: verifiable against firmware_app/src/lib.rs and
%  firmware_app/src/residual_nn.rs. Where this chapter states a constant or an equation,
%  it is the one in the code, not an idealisation of it. If the code changes, this
%  chapter changes.
%
%  FRAME CONVENTION: z-up (ENU), gravity acts along -e_3. This matches the implementation
%  and Crazyswarm2. Note that docs/04_Unified_Residual_Wrench_Model.md was written with a
%  +g e_3 term; that document predates this chapter and is superseded on this point.
%
%  Citation baseline: 2026-08-27 (see docs/thesis/CITATION_BASELINE.md).
%  After that date: only NEW cites or EDITED prior-work sentences need re-checking.
% =====================================================================================

\chapter{System Modelling and Residual-Force Formulation}
\label{ch:model}

Chapter~\ref{ch:related} described methods that either measure or predict the aerodynamic
interaction between vehicles. This chapter defines precisely what is being measured or predicted,
so that the strategies compared in Chapter~\ref{ch:control} can be said to estimate the
\emph{same} quantity. That is a prerequisite for the comparison: two controllers that compensate
differently-defined residuals are not comparable, however similar their block diagrams look.

The chapter also states the assumptions the formulation rests on, and where each one is expected
to fail. Several are known to be imperfect --- the summation in
Section~\ref{sec:model-deepsets} most conspicuously --- and stating them here is what allows
the discussion chapter to interpret a negative result as evidence about an assumption
rather than as an unexplained failure.

% -------------------------------------------------------------------------------------
\section{Notation and conventions}
\label{sec:model-notation}

Vectors are expressed in a world frame that is \emph{east-north-up}: the third basis vector $e_3$
points away from the ground and gravity acts along $-e_3$. This matches the convention used by the
flight software and the motion-capture system, and is stated explicitly because the opposite
convention is common in the quadrotor literature and the sign of the gravity term differs between
them.

For vehicle $i$: $p_i \in \mathbb{R}^3$ is position, $v_i \in \mathbb{R}^3$ velocity,
$R_i \in SO(3)$ the rotation from body to world, and $\omega_i \in \mathbb{R}^3$ the body angular
velocity. The hat map $\hat{\cdot} : \mathbb{R}^3 \to \mathfrak{so}(3)$ is defined by
$\hat{a}b = a \times b$. The vehicles are homogeneous, of mass $m$ and diagonal inertia $J$.
Rotor $k$ spins at $\varpi_k$ and produces thrust $k_{t,k}\varpi_k^2$ along the body $z$-axis.

% -------------------------------------------------------------------------------------
\section{Rigid-body model with a residual wrench}
\label{sec:model-dynamics}

Each vehicle is modelled as a rigid body driven by a collective thrust $f_i$ along its body
$z$-axis and a body torque $\tau_i$:
%
\begin{align}
  \dot{p}_i &= v_i, &
  m\,\dot{v}_i &= -m g e_3 + f_i R_i e_3 + f_{\mathrm{res},i}, \label{eq:model-trans}\\
  \dot{R}_i &= R_i \hat{\omega}_i, &
  J\dot{\omega}_i &= -\omega_i \times J\omega_i + \tau_i + \tau_{\mathrm{res},i}.
  \label{eq:model-rot}
\end{align}
%
The terms $f_{\mathrm{res},i}$ and $\tau_{\mathrm{res},i}$ form the \emph{residual wrench}: by
construction they collect every force and moment that \eqref{eq:model-trans}--\eqref{eq:model-rot}
otherwise omits. This includes blade flapping, rotor drag, ground effect and errors in $m$, $J$ and
$k_t$, as well as the inter-vehicle aerodynamic interaction that is the subject of this thesis.

That the residual is a catch-all rather than a physical quantity has two consequences that matter
for the experiments. First, a residual measured in isolated free flight is not zero; it is whatever
the nominal model gets wrong about a single vehicle. Second, the interaction can only be isolated
by \emph{differencing} --- comparing the residual measured in formation against the residual
measured alone under the same trajectory. The coplanar control scenarios in
Chapter~\ref{ch:setup} exist for exactly this purpose.

\subsection{Scope: force, not moment}

This thesis compensates the \emph{force} residual $f_{\mathrm{res}}$. The moment residual
$\tau_{\mathrm{res}}$ is measured and logged, but no strategy compensates it separately: the
attitude loop rejects it as a disturbance in the usual way. This is a scoping decision, not a claim
that $\tau_{\mathrm{res}}$ is negligible --- the characterisation literature measures interaction
\emph{moments} alongside forces on a test rig~\cite{kiran2025downwash}, and the discussion chapter
returns to what leaving them uncompensated may cost.

% -------------------------------------------------------------------------------------
\section{Measuring the residual}
\label{sec:model-measure}

The residual can be recovered without any model of the interaction, by comparing what the vehicle
actually accelerated with what the nominal model says it should have accelerated given the
actuation actually applied. Rearranging \eqref{eq:model-trans},
%
\begin{equation}
  \frac{f_{\mathrm{res},i}}{m}
  \;=\; \underbrace{\dot{v}_i}_{a_{\mathrm{meas}}}
      - \underbrace{\left( \frac{f_i}{m} R_i e_3 - g e_3 \right)}_{a_{\mathrm{model}}} .
  \label{eq:model-ares}
\end{equation}
%
Both terms are available onboard. The measured acceleration comes from the accelerometer, which
senses specific force in the body frame and therefore excludes gravity by construction; rotating it
into the world frame and subtracting $g e_3$ gives $a_{\mathrm{meas}}$. The model acceleration
requires the applied thrust, which is reconstructed from measured rotor speeds as
$f_i = \sum_{k=1}^{4} k_{t,k}\varpi_k^2$.

The implementation evaluates \eqref{eq:model-ares} at every control step and logs it as
$a_{\mathrm{res}}$, in every controller mode --- including modes that make no use of it. This is
deliberate. The training data for every predictive strategy is collected under the
\emph{uncompensated} geometric controller, so the measurement must exist when no INDI term is
active; and the comparison of a predicted residual against a measured one, which is how a learned
model is evaluated, requires both quantities on the same flight.

\subsection{What this measurement requires, and how it fails}
\label{sec:model-measure-fail}

Three properties of \eqref{eq:model-ares} shape the experimental design.

\paragraph{It requires rotor-speed sensing.} Without $\varpi_k$ the applied thrust is unknown and
$a_{\mathrm{model}}$ cannot be formed. This is the hardware prerequisite identified in
Section~\ref{sec:rw-reactive}, and the same one that Cobo-Briesewitz et
al.~\cite{cobobriesewitz2026lindi} remove by learning the residual instead of computing it. In the
implementation, the absence of a rotor-speed source causes $a_{\mathrm{res}}$ to read
\emph{identically zero}. That failure is silent and is indistinguishable, in a log file, from a
flight in which no interaction occurred. Verifying that $a_{\mathrm{res}} \neq 0$ is therefore the
first check of every experimental session (Chapter~\ref{ch:setup}).

\paragraph{It is delayed.} $a_{\mathrm{meas}}$ is differentiated from noisy inertial measurement
and must be filtered; the reconstructed thrust inherits the rotor-speed sensing delay. The estimate
therefore exists only after the disturbance has begun to act. No amount of tuning removes this: it
is the structural property that motivates prediction.

\paragraph{It absorbs model error.} Any error in $m$ or $k_t$ appears in $a_{\mathrm{model}}$ and is
attributed to $f_{\mathrm{res}}$. A miscalibrated thrust constant produces a residual that looks
like a constant interaction force. This is why the airframe constants are identified once and
frozen, and why the null scenarios matter: a residual that persists when no neighbour is present is
model error, not interaction.

\subsection{Sign convention in the control law}
\label{sec:model-sign}

The compensation enters with a sign that is easy to get wrong and whose error is not obvious in
flight. From \eqref{eq:model-trans}, holding a desired acceleration $a_{\mathrm{des}}$ requires
%
\begin{equation}
  f_i R_i e_3 \;=\; m\,a_{\mathrm{des}} + m g e_3 - f_{\mathrm{res},i},
  \label{eq:model-sign}
\end{equation}
%
so the desired-acceleration vector passed to the thrust-and-attitude allocation carries
$-f_{\mathrm{res},i}/m$. The residual is \emph{subtracted}, whether it was measured
(Section~\ref{sec:model-measure}) or predicted (Section~\ref{sec:model-deepsets}).

Adding it instead does not merely fail to compensate: it doubles the disturbance, because the
controller then commands a thrust error equal and opposite to the correction it should have made.
The resulting behaviour is a stable, well-controlled vehicle that tracks slightly worse in
formation --- which is what an uncompensated vehicle also looks like. In single-vehicle flight the
error is invisible altogether, since $f_{\mathrm{res}} \approx 0$ with no neighbour present. This is
recorded here because it is a property of the formulation rather than a coding detail, and because
it establishes a specific prediction that the two-robot results chapter can check: with the sign
correct, position-INDI compensation must \emph{reduce} the formation-induced displacement, and with
it inverted the displacement must be approximately twice the uncompensated value.

% -------------------------------------------------------------------------------------
\section{Predicting the residual from relative state}
\label{sec:model-deepsets}

The interaction force on vehicle $i$ is a function of the configuration of the team relative to it,
%
\begin{equation}
  f_{\mathrm{res},i} \;\approx\; \mathcal{F}\!\left( \{\, p_j - p_i,\; v_j - v_i \,\}_{j \neq i} \right),
  \label{eq:model-F}
\end{equation}
%
and \eqref{eq:model-F} is what a predictive strategy learns. Two structural properties are required
of any parameterisation of $\mathcal{F}$, and both follow from the physics rather than from
convenience. The value must not depend on the order in which neighbours are enumerated, and it must
be defined for a varying number of them, since a vehicle in a three-robot team has two neighbours
and one in a pair has one.

\subsection{Deep-sets parameterisation}

The parameterisation used here satisfies both by construction:
%
\begin{equation}
  \hat{a}_{\mathrm{res},i}
  \;=\; \rho\!\left( \sum_{j \neq i} \phi\!\left( p_j - p_i,\; v_j - v_i \right) \right),
  \label{eq:model-deepsets}
\end{equation}
%
where $\phi$ is applied to each neighbour independently and $\rho$ maps the summed embedding to a
residual acceleration in the world frame. Summation is permutation-invariant and accepts any number
of terms, so one parameter set serves teams of different sizes. The output is an acceleration rather
than a force so that it is directly comparable with the measured $a_{\mathrm{res}}$ of
\eqref{eq:model-ares}; the two quantities are in the same units, in the same frame, at the same
point in the control loop.

The concrete network is small: $\phi$ maps $6 \rightarrow 16 \rightarrow 16 \rightarrow 8$ and
$\rho$ maps $8 \rightarrow 16 \rightarrow 16 \rightarrow 3$, with rectified-linear activations
throughout except at the output, which is linear because an interaction force may point either way.
This is \num{987} parameters, evaluated onboard within the control period. The sizing is discussed
in Chapter~\ref{ch:control}; what matters here is that \eqref{eq:model-deepsets} is the model,
and the layer widths are an implementation choice within it.

\paragraph{These equations are the code.} Equation~\eqref{eq:model-deepsets}, the layer widths
above, the distance guards of Section~\ref{sec:model-guards} and the output bound are the values
in \texttt{firmware\_app/src/residual\_nn.rs}, not an idealisation of them, and were read from it
rather than reconstructed. The same applies to \eqref{eq:model-ares} and the sign convention of
Section~\ref{sec:model-sign}, which correspond to the residual computation and the
desired-acceleration assembly in \texttt{firmware\_app/src/lib.rs}. If the implementation changes,
this chapter changes with it; a modelling chapter that has drifted from the flight code describes a
system that was never flown.

\subsection{The superposition assumption}
\label{sec:model-superposition}

Equation~\eqref{eq:model-deepsets} assumes the aggregate interaction is a sum of per-neighbour
contributions in the latent space of $\phi$. This is more general than summing forces directly ---
$\rho$ is free to be nonlinear in the sum --- but it is still an assumption, and the literature
states plainly that it does not hold exactly. Shi et al.~\cite{shi2020neuralswarm} report that with
more than two vehicles the aerodynamic effect is not a simple superposition of each pair, and Gielis
et al.~\cite{gielis2023aggregate} motivate a learned aggregation model, measured on a load stand,
precisely on the grounds that a single-vehicle downwash model is unlikely to suffice for predicting
aggregate forces.

The assumption is retained here for two reasons. It is what makes one trained model applicable to
teams of two and three without retraining, which the comparison requires. And a controller does not
need an exact model: it needs one whose residual after compensation is smaller than the residual
before it. The relevant question is therefore not whether \eqref{eq:model-deepsets} is exact, which
it is not, but how large the resulting error is on a deployed model --- which is what
the multi-robot results chapter measures.

\subsection{Input conditioning}
\label{sec:model-guards}

Two guards are applied to the inputs of $\phi$, and both are part of the model rather than
defensive programming, because they change the function being learned.

A neighbour farther than \SI{2}{\metre} is excluded from the sum. Beyond that distance the
interaction is below the noise floor of the measurement, and including such neighbours would ask
the network to learn that distant vehicles do nothing --- spending capacity on a region where the
answer is known.

A separation closer than \SI{4}{\centi\metre} is clamped to that value, preserving direction. Below
it the model would extrapolate beyond anything the training data can contain, since the vehicles are
\SI{9}{\centi\metre} across and such a configuration is a collision. A learned function is least
trustworthy exactly where the physics is strongest, and the clamp bounds the extrapolation rather
than trusting it.

The output is bounded as well, at \SI{8}{\metre\per\second\squared} --- comfortably above the
interaction accelerations expected between vehicles of this size, and far below what an
ill-conditioned network can emit. This is not a hypothetical failure mode. In simulation, Gu et
al.~\cite{gu2025comparative} report a radial-basis-function estimator crashing the vehicle under
periodic wind, and an extreme learning machine becoming unstable on landing under ground effect;
their echo state network and Gaussian process did not. The authors say the failure \emph{may be
attributed} to the random initialisation of the input-to-hidden weights, which \emph{``can produce
excessively large prediction values''}. Their vehicle is a single quadrotor and their disturbances
are wind and ground effect, so this is not evidence about inter-vehicle interaction --- but it is
evidence that an unbounded learned term can destabilise a controller. A saturating output converts
that failure into a bounded and, because saturation is logged, a visible one.

The training data must pass through the same guards as the deployed model. A model fitted to
unguarded inputs is fitted to a different distribution than the one it will meet, and the
discrepancy is invisible until predicted and measured residuals are compared in flight.

\subsection{Normalisation}

Input scaling is folded into the first layer of $\phi$ at export rather than applied onboard. For
weights $W_1$ and bias $b_1$ trained on inputs normalised by mean $\mu$ and standard deviation
$\sigma$,
%
\begin{equation}
  W_1 \frac{x - \mu}{\sigma} + b_1
  \;=\; \underbrace{\left(W_1 \oslash \sigma\right)}_{W_1'} x
      + \underbrace{\left(b_1 - \left(W_1 \oslash \sigma\right)\mu\right)}_{b_1'},
  \label{eq:model-fold}
\end{equation}
%
with $\oslash$ denoting column-wise division. The two forms are arithmetically identical, so the
deployed network computes the same function as the trained one while the flight software carries no
normalisation constants that could fall out of step with the model.

% -------------------------------------------------------------------------------------
\section{Assumptions and where they fail}
\label{sec:model-assumptions}

Table~\ref{tab:model-assumptions} collects the assumptions this formulation makes. Each is stated
with the condition under which it is expected to break, so that an experimental result can be
attributed to a specific assumption rather than to the model as a whole.

\begin{table}[t]
  \centering
  \caption{Assumptions of the residual formulation and their expected failure modes.}
  \label{tab:model-assumptions}
  \small
  \begin{tabular}{@{}p{0.30\linewidth}p{0.30\linewidth}p{0.32\linewidth}@{}}
    \toprule
    Assumption & Where it fails & Consequence \\
    \midrule
    Rigid body; no aeroelastic deformation
      & Not expected to fail at this scale
      & --- \\
    $m$, $J$, $k_t$ known and constant
      & Battery discharge; rotor wear; miscalibration
      & Model error appears as a constant residual; bounded by the null scenarios \\
    Residual is dominated by inter-vehicle interaction
      & Near the ground; near walls
      & Flight volume is bounded away from both (Chapter~\ref{ch:setup}) \\
    Neighbour relative states are known accurately
      & Motion-capture dropout; latency
      & Prediction degrades; measurement is unaffected, which is itself a comparison result \\
    Aggregate interaction superposes in latent space
      & Three or more vehicles, overlapping wakes
      & \textbf{Known to be imperfect}~\cite{shi2020neuralswarm,gielis2023aggregate}; magnitude measured in the multi-robot results chapter \\
    Interaction lies within the training distribution
      & Unseen separations, faster relative motion
      & Prediction degrades while measurement does not; this is the crossover of \textbf{RQ2} \\
    Rotor-speed measurement available
      & No RPM source fitted
      & $a_{\mathrm{res}}$ reads identically zero --- silent, and checked every session \\
    \bottomrule
  \end{tabular}
\end{table}

Two entries deserve emphasis because they are not merely caveats but hypotheses the experiments
test. The superposition assumption is known to be imperfect, and the experimental question is the
size of the penalty rather than its existence. And the training-distribution assumption is exactly
what separates the predictive family from the reactive one: a measurement has no training
distribution to leave. Cobo-Briesewitz et al.~\cite{cobobriesewitz2026lindi} observe precisely this
effect for their purely learned method in the no-payload case, attributing a drop in performance on
the circular trajectory to its being out of distribution relative to the training data.

% -------------------------------------------------------------------------------------
\section{Summary}
\label{sec:model-summary}

The residual wrench of \eqref{eq:model-trans}--\eqref{eq:model-rot} is the single quantity every
strategy in this thesis estimates. It can be recovered online from
\eqref{eq:model-ares}, which needs no model of the interaction but requires rotor-speed sensing and
is available only after the disturbance acts; or predicted from \eqref{eq:model-deepsets}, which
acts ahead of the disturbance but is valid only within its training distribution and assumes an
aggregation that is known to be approximate. Both are expressed as accelerations in the world frame
and enter the control law at the same point with the same sign, \eqref{eq:model-sign}.

That common definition is what makes the comparison in the two-robot results chapter a
comparison rather than a collection of separately-tuned controllers. Chapter~\ref{ch:control}
describes the seven strategies built on it, and the point at which each injects its knowledge of
the interaction.
